% sec-04: the robotic Phase 1 in numbers, and the Specific Aims in brief.
% Table 24.
\section{The Robotic Phase 1, in Numbers}

Every letter of the week points back to one trial, and a funder deciding whether
it is worth money needs to see that every part of it is specified to a number
someone can check. Table 24 gathers those numbers from the deposited protocol,
the investigational new drug application, the capitalization plan and the site
package \cite{onpremwhippl,phase1ind,capitalplan2026,movein2026}. Nothing in it
is new; everything in it is checkable.

\subsection{What the numbers say to a funder}

The design is conventional where it should be and specified where it is new. A
3+3 escalation across three dose levels in up to 18 participants is the most
familiar Phase 1 design there is. What is new, the robot and the model, is bound
by numbers: a 3 ms stop across eight arms, a tip-force cap of 3 N per arm and
18 N in all, a five-vessel no-fly gate, a gate of at least 1,000 simulated
procedures across at least two independent frameworks with a Unified Safety
Level of at least 7.0 before any patient, and a language model with no route to
act \cite{onpremwhippl,phase1ind}. A funder does not have to trust the
company's judgment about safety; it can read the limits.

The evidence behind the agent is a dated chronology. In August 2025 the
company's ten-arm simulation returned 12.8 against 5.4 months of median overall
survival \cite{qspmetpancre}; in May 2026 the developer's Phase 3 trial reported
13.2 against 6.6 in the RAS G12 population \cite{rasolute302}; and the approval
that followed reported 13.2 against 6.7 in the overall population
\cite{fdarasonque2026}. The company calls the closeness a chronology, not a
validation, because the populations, regimens and sample sizes differ.
Pancreatic ductal adenocarcinoma remains the cancer with the widest gap between
incidence and survival in the United States \cite{siegel2025statistics}, and
after resection the standard remains combination chemotherapy
\cite{conroy2018folfirinox}.

\begin{apptable}
\begin{tabularx}{\textwidth}{>{\raggedright\arraybackslash}p{5.0cm} >{\raggedright\arraybackslash}p{7.3cm} >{\raggedright\arraybackslash}X}
\headrow \thead{Parameter} & \thead{Value} & \thead{Source}\\
\midrule
Design & Open-label, single-arm, 3+3 escalation & Protocol \\
Dose levels of daraxonrasib & 160, 220 and 300 mg & Protocol \\
Participants & Up to 18, with sentinel enrollment & Protocol \\
Dose-limiting toxicity window & 28 days & Protocol \\
Primary safety window & Day 30; Clavien-Dindo III or higher & Protocol \\
Key secondary endpoints & R0 at day 90; fistula; 24-month survival & Protocol \\
\midrule
Arms and degrees of freedom & 8 arms, 56 degrees of freedom & Protocol \\
Sensor channels & 640, at 80 per arm & Protocol \\
Positional accuracy & 0.05 mm root mean square at 1,200 mm/s & Protocol \\
Emergency stops & 3 ms across arms; 500 ms for the system & Protocol, IND \\
Tip-force caps & 3 N per arm; 18 N across all arms & Protocol \\
Gate before any patient & 1,000 simulated procedures, 2 frameworks & IND \\
Language model & Advisory only; no write access, no route & IND \\
\midrule
Simulated survival, 2025 & 12.8 against 5.4 months; hazard ratio 0.25 & Company QSP \\
Credibility battery & 81.9 over 55 tests & Capital plan \\
Reported survival, RAS G12 & 13.2 against 6.6 months & RASolute 302 \\
\midrule
NIH Phase I, then Phase II & \$306,000 over 9 months, then \$1,300,000 & Capital plan \\
Clinical conduct in Phase II & 6 participants at \$48,000 each & Capital plan \\
Program direct cost & \$700,000 a year, \$521,000 of it personnel & Site package \\
\end{tabularx}
\end{apptable}
\vspace{-0.60cm}
\tabcap{Table 24. The robotic Phase 1 in nineteen numbers, grouped as design, procedure and\\
safety, evidence, and money, each with the deposited document a funder can check.}

\subsection{The Specific Aims, in brief}

The one-page Specific Aims sent with today's letter 1 proposes three aims for an
NIH Phase I of \$306,000 over nine months \cite{capitalplan2026}. The first fixes
the model's context of use inside the protocol: the decisions it may advise on,
and those it may never touch, which are dose assignment, dose-limiting toxicity
calls and analysis. The second evaluates the layer in simulated use with the
site's surgical team, as part of the protocol's simulation gate. The third
assembles the credibility evidence for that context of use, framed on ASME V\&V
40 \cite{asmevv40} and read against the FDA's draft guidance on artificial
intelligence in regulatory decision-making \cite{fdaaiguidance2025}, citing
rather than repeating any NSF-funded verification results \cite{nsf26510}. None
of the three aims asks NIH to pay for what the NSF Pitch proposes.
